\documentclass[12pt,a4paper,oneside]{article}
\usepackage[brazil]{babel}
\usepackage[utf8]{inputenc}
\usepackage{olymp}

\usepackage{wrapfig}

\setcounter{problem}{0}

\begin{document}

\begin{problem}{Engrenagens}{engrenagens}{2}

\begin{wrapfigure}{r}{0.3\textwidth}
    \includegraphics[width=0.3\textwidth]{images/exercise_7_0}
\end{wrapfigure}

O setor de teste de uma fábrica de engrenagens recebe diariamente centenas de amostras de engrenagens de vários tipos para verificação, sendo que, para cada tipo, existem vários testes a serem realizados. Um dos tipos de engrenagens testado é a engrenagem conjugada, em que duas engrenagens são fixadas uma a outra, com o objetivo de transferir torque/velocidade. Existem inúmeras combinações de tamanhos de engrenagens, e um requisito imposto pelo cliente é que, quando a engrenagem maior der uma volta completa, a menor tenha dado um número inteiro de voltas completas. A engrenagem conjugada ao lado segue este requisito, pois quando a maior (de 21 dentes) terminar uma volta, a menor (de 7 dentes) terá dado exatamente 3 voltas.

Para automatizar os testes, a fábrica já tem uma máquina que conta a quantidade de dentes de cada engrenagem. Agora precisa de um programa que verifique se atendem ao requisito. Note que, para que isto aconteça, a quantidade de dentes da maior deve ser um múltiplo da quantidade de dentes da menor, ou seja, a divisão dessas quantidades deve ser exata. 


%\Notes
%
%Observações, se tiver.

\InputFile

A entrada é composta por dois valores inteiros, $D_1$ e $D_2$, que são a quantidade de dentes respectivamente da menor e da maior engrenagem. Restrição: $0 ≤\leq D_1, D_2 \leq 5000$.

%\Notes


\OutputFile

Seu programa deve gerar apenas uma linha de saída, contendo um valor inteiro, sendo `\texttt{1}' para aceitar a engrenagem e `\texttt{0}' para rejeitá-la.

\Example

\begin{example}
\exmp{
7 21
}{
1
}%
\end{example}


\begin{example}
\exmp{
5 50
}{
1
}%
\end{example}

\begin{example}
\exmp{
12 48
}{
1
}%
\end{example}

\begin{example}
\exmp{
20 50
}{
0
}%
\end{example}

\end{problem}

\end{document}
